% Drawing 119/08-6 -- folder 119, batch 08, archive page 156 (PDF page 157).
% Transcription: carré à deux diagonales, avec $\varepsilon^*$, $a_Q$, $Z_Q$
% Kind: square with two diagonals. Method 2 (schematic redraw from the facsimile). Status: ok.
% Notes: Between parentheses, after a colon: a small square with both diagonals; epsilon*, a_Q and Z_Q below.
% Coordinates: 1 unit = 1 mm = 1/100 of the crop width (bbox in page fractions: [0.41, 0.19, 0.54, 0.27]).
% Redrawn by Claude (claude-opus-5-5), issue #73, 2026-11-07; not checked by a human; not compiled (no TeX in the session).
% Rights: reproduces part of the work -- behind RENDER_TRANSCRIPTS, never published by itself.
\documentclass[tikz,border=3pt]{standalone}
\usepackage{amsmath,amssymb}
\usetikzlibrary{patterns,arrows.meta}
\tikzset{pencil/.style={black!45}, struck/.style={black!35}, mark/.style={circle,draw,inner sep=1.3pt}, pt/.style={circle,fill,inner sep=0.9pt}, opt/.style={circle,draw,fill=white,inner sep=1.1pt}}
\begin{document}
\begin{tikzpicture}[x=1mm,y=1mm,line cap=round,line join=round,font=\scriptsize]
\draw plot[smooth, tension=0.55] coordinates {(10.8,-19.9) (4.6,-35.3) (10.8,-50.7)};
\node[inner sep=1pt] at (20,-29.1) {\tiny :};
\draw[very thick] (38.5,-23) -- (60,-23) -- (60,-38.4) -- (38.5,-38.4) -- cycle;
\draw (38.5,-23) -- (60,-38.4);
\draw (38.5,-38.4) -- (60,-23);
\draw plot[smooth, tension=0.55] coordinates {(93.8,-19.9) (100,-35.3) (93.8,-53.7)};
\node[inner sep=1pt] at (20,-50.7) {\tiny $\varepsilon^*$};
\node[inner sep=1pt] at (47.7,-63) {\tiny $a_Q$};
\node[inner sep=1pt] at (75.4,-59.9) {\tiny $Z_Q$};
\end{tikzpicture}
\end{document}
